% Generated draft: requires mathematical and source review.
\documentclass[11pt]{article}
\usepackage[margin=25mm]{geometry}
\usepackage{fontspec}
\setmainfont{Latin Modern Roman}
\usepackage{amsmath,amssymb,mathtools,unicode-math}
\setmathfont{Latin Modern Math}
\usepackage{enumitem,needspace,xurl,hyperref}
\setlength{\parindent}{0pt}
\setlength{\parskip}{5pt}
\setlength{\emergencystretch}{3em}
\newcommand{\N}{\mathbb{N}}
\newcommand{\Z}{\mathbb{Z}}
\newcommand{\Q}{\mathbb{Q}}
\newcommand{\R}{\mathbb{R}}
\newcommand{\C}{\mathbb{C}}
\newcommand{\F}{\mathbb{F}}
\newcommand{\A}{\mathbb{A}}
\newcommand{\PP}{\mathbb{P}}
\newcommand{\E}{\mathbb{E}}
\newcommand{\eps}{\varepsilon}
\newcommand{\dd}{\,\mathrm{d}}
\begin{document}
\section*{UnsolvedMath ID: 800788}
\textbf{Source problem code:} catalogue entry
\subsection{Definitions and mathematical statement}
Let \(n\geq 1\), let \(O\) denote the origin of \(\C^n\), and let \(\operatorname{End}(\C^n,O)\) denote the holomorphic map germs \(f:(\C^n,O)\to(\C^n,O)\) fixing \(O\). A germ is dominant if its complex Jacobian determinant is not identically zero, and is non-invertible if it has no holomorphic inverse germ; equivalently, \(\det Df(O)=0\). Write \(f^k\) for the \(k\)-fold iterate of \(f\).

The supplied excerpt does not define \(c\). Under the standard order-of-vanishing interpretation, for a nonzero holomorphic germ \(g=(g_1,\ldots,g_n)\) fixing \(O\), set
\[
 c(g)=\min_{1\leq i\leq n}\operatorname{ord}_O(g_i),
 \qquad
 \operatorname{ord}_O(h)=\min\{\alpha_1+\cdots+\alpha_n:a_\alpha\neq0\},
\]
where \(h(z)=\sum_{\alpha\in\Z_{\geq0}^n}a_\alpha z_1^{\alpha_1}\cdots z_n^{\alpha_n}\), with \(\operatorname{ord}_O(0)=+\infty\). Dominance ensures that this order is finite for every iterate.

With this interpretation of \(c\), Abate's Problem OP6 asks: for every non-invertible dominant germ \(f\in\operatorname{End}(\C^n,O)\), does the sequence \(\{c(f^k)\}_{k\geq1}\) eventually satisfy a linear recurrence with integer coefficients? In the usual homogeneous, constant-coefficient, monic convention, this means that there exist integers \(r\geq1\), \(K\geq1\), and \(b_0,\ldots,b_{r-1}\in\Z\), which may depend on \(f\), such that
\[
 \forall k\geq K,\qquad
 c(f^{k+r})=\sum_{j=0}^{r-1}b_j\,c(f^{k+j}).
\]
The definition of \(c\) and the precise recurrence convention require confirmation against the source: neither is specified in the supplied excerpt.
\subsection{Short English statement}
Abate asks whether the orders of vanishing of the iterates of every non-invertible dominant holomorphic germ fixing the origin eventually obey a fixed linear recurrence with integer coefficients. This formulation uses an explicitly provisional interpretation of the undefined notation \(c\) and of the recurrence convention.
\subsection{Sources}
\begin{itemize}
\item Abate, \emph{Open Problems in Local Discrete Holomorphic Dynamics} (2011), Problem OP6. Supplied source URL: \url{https://arxiv.org/abs/1107.4876}. The supplied provenance records access on 2026-10-16 and extraction by \texttt{arxiv\_latex}; these details are not independently verified here.
\item Local UnsolvedMath registry, AMR Open Problem Lists: permanent ID \(800788\), source problem code \texttt{catalogue entry}, upstream ID \(800788\), upstream revision \texttt{d371ee917ac9647fbb06cd6fc529deec6f3dda82}. The record attributes the problem to Abate (2011) and supplies the status \texttt{open}, with record update date 2026-10-16. This is a supplied registry label, not an independently verified assessment of the literature.
\end{itemize}
\end{document}
